Under the binary recoding in which Saco's treatment indicator is one for the present arm \(1\) and zero for the present arm \(2\), the bounded two-regime experiment constructed in the preceding proposition satisfies every assumption explicitly listed in Saco's Theorem 5.14 and also its maintained consistency and no-interference condition. Indeed, take the deterministic scored set to be all \(T\) occasions, take both predictable nuisance regressions to be identically zero, take \(Y_t(0)=Y_t(1)=1\) in Saco's binary notation, and take his target to be \(\theta_0=0=\theta^*\). His centered AIPW increment is then exactly \(\xi_t\); conditionally on \(\mathcal F_{t-1}\), the observable centered score has the nondegenerate two-point law \(P(\xi_t=1/p_t\mid\mathcal F_{t-1})=p_t\) and \(P(\xi_t=-1/(1-p_t)\mid\mathcal F_{t-1})=1-p_t\); hence \(E(\xi_t\mid\mathcal F_{t-1})=0\) and \(\operatorname{Var}(\xi_t\mid\mathcal F_{t-1})=1/p_t+1/(1-p_t)\in\{4,16/3\}\), and his realized and predictable quadratic variations are \(Q_T=[S]_T\) and \(V_T^2=\langle S\rangle_T\), with \(V_T^2\ge 4T\) almost surely. Nevertheless, along \(T=2n\),
\[
 \frac{S_T}{\sqrt{Q_T}}\Longrightarrow
 \mathbf1\{Z>0\}\frac{2Z+4W/\sqrt3}{\sqrt{28/3}}
 +\mathbf1\{Z\le0\}\frac{2Z+2W}{\sqrt8},
\]
where \(Z,W\) are independent \(N(0,1)\) variables and the displayed limit has nonzero mean. The feasible sample-variance Studentized AIPW statistic has the same limit. Therefore the variance-growth-only Studentized central limit conclusion asserted in Saco's Theorem 5.14 is false as stated.